%% For double-blind review submission, w/o CCS and ACM Reference (max submission space)
%\documentclass[acmsmall,review]{acmart}\settopmatter{printfolios=true,printccs=false,printacmref=false}
%% For double-blind review submission, w/ CCS and ACM Reference
\documentclass[acmsmall,review,anonymous,screen]{acmart}\settopmatter{printfolios=true,printacmref=false}
%% For double-blind review submission, w/ CCS and ACM Reference
%\documentclass[acmsmall,screen]{acmart}\settopmatter{printfolios=true,printacmref=false}
%% For single-blind review submission, w/o CCS and ACM Reference (max submission space)
%\documentclass[acmsmall,review]{acmart}\settopmatter{printfolios=true,printccs=false,printacmref=false}
%% For single-blind review submission, w/ CCS and ACM Reference
%\documentclass[acmsmall,review]{acmart}\settopmatter{printfolios=true}
%% For final camera-ready submission, w/ required CCS and ACM Reference
%\documentclass[acmsmall]{acmart}\settopmatter{}
\usepackage[shortlabels]{enumitem}
\usepackage{mathtools}
\usepackage{wrapfig}
\usepackage{stackengine}
\usepackage{ textcomp }

\DeclareSymbolFont{arrows}{U}{FdSymbolC}{m}{n}

% \usepackage{ amssymb }
\DeclareTextFontCommand{\texttt}{\ttfamily}

%% Journal information
%% Supplied to authors by publisher for camera-ready submission;
%% use defaults for review submission.
\acmJournal{PACMPL}
\acmVolume{}
\acmNumber{OOPSLA} % CONF = POPL or ICFP or OOPSLA
\acmArticle{}
\acmYear{2025}
\acmMonth{1}
\acmDOI{} % \acmDOI{10.1145/nnnnnnn.nnnnnnn}
\startPage{1}

%% Copyright information
%% Supplied to authors (based on authors' rights management selection;
%% see authors.acm.org) by publisher for camera-ready submission;
%% use 'none' for review submission.
\setcopyright{none}
%\setcopyright{acmcopyright}
%\setcopyright{acmlicensed}
%\setcopyright{rightsretained}
%\copyrightyear{2018}           %% If different from \acmYear

%% Bibliography style
\bibliographystyle{ACM-Reference-Format}
%% Citation style
%% Note: author/year citations are required for papers published as an
%% issue of PACMPL.
%\citestyle{acmauthoryear}   %% For author/year citations
\citestyle{acmnumeric}

%%%%%%%%%%%%%%%%%%%%%%%%%%%%%%%%%%%%%%%%%%%%%%%%%%%%%%%%%%%%%%%%%%%%%%
%% Note: Authors migrating a paper from PACMPL format to traditional
%% SIGPLAN proceedings format must update the '\documentclass' and
%% topmatter commands above; see 'acmart-sigplanproc-template.tex'.
%%%%%%%%%%%%%%%%%%%%%%%%%%%%%%%%%%%%%%%%%%%%%%%%%%%%%%%%%%%%%%%%%%%%%%


%% Some recommended packages.
\usepackage{booktabs}   %% For formal tables:
                        %% http://ctan.org/pkg/booktabs
\usepackage{subcaption} %% For complex figures with subfigures/subcaptions
                        %% http://ctan.org/pkg/subcaption
 \usepackage{ stmaryrd }                       

\usepackage{relsize}
\usepackage{mathpartir}
\usepackage{amsmath}
\usepackage{amsthm}
\usepackage{listings}
\usepackage{xspace}
\usepackage{definitions}
\usepackage{multirow,bigdelim}
\usepackage{pbox}
\usepackage{courier}
\usepackage{soul}
\usepackage{centernot}
\usepackage{dutchcal}
%\usepackage{pzccal} 
%\DeclareFontFamily{OT1}{pzc}{}
%\DeclareFontShape{OT1}{pzc}{m}{it}{<-> s * [1.10] pzcmi7t}{}
%\DeclareMathAlphabet{\mathpzc}{OT1}{pzc}{m}{it}

\newcommand\multibrace[3]{\rdelim\}{#1}{3mm}[\pbox{#2}{#3}]}


%%these COLOUR MACROS ARE ACTIVELY FUCKING EVIL
%%DO NOT DO THIS. EVER
%%AT LEast have some ovious way to TURN THEM OFF

\definecolor{ferngreen}{rgb}{0.31, 0.47, 0.26}
\newcommand{\kjx}[1]{{}}
\newcommand{\scd}[1]{{}}
%\newcommand{\sdN}[1]{{\color{dkgreen}{#1}}}
%\newcommand{\jm}[1]{{\color{magenta}{JM: #1}}}
\newcommand{\sdcomment}[1]{{}}
\newcommand{\secomment}[1]{{}}
\newcommand{\jncomment}[1]{{}}

\newcommand{\sd}[1]{{{#1}}}
\newcommand{\se}[1] {#1} %{{\color{brown}{#1}}}
\newcommand{\sue}[1] {{\color{brown}{#1}}}
\definecolor{amber}{rgb}{1.0, 0.75, 0.0}
\definecolor{amethyst}{rgb}{0.6, 0.4, 0.8}

\newcommand{\ponders}[3]{\marginpar{\tiny\itshape\raggedright\textcolor{#2}{\textbf{#1:} #3}}\ignorespaces}
\renewcommand{\ponders}[3]{{#3\ignorespaces}}  %TURNING IT OFF. NOT THAT IT FUCKING HEPLS

\marginparwidth=1.6cm \marginparsep=0cm
\newcommand{\TODO}[1]{} % {{\color{red}#1}}
\newcommand{\sophia}[1]{{#1}}
\newcommand{\toby}[1]{} % {\ponders{Toby}{purple}{#1}}
\newcommand{\susan}[2][]{{#2}\xspace}
\newcommand{\james}[1]{\ponders{James}{orange}{#1}}
\newcommand{\jm}[2][]{\ponders{Julian}{magenta}{#1} {#2}\xspace}
\newcommand{\mrr}[2][]{\ponders{Matthew Ross}{offblue}{{#1}} {{#2}}\xspace}
\newcommand{\sdNr}[2][]{\ponders{SD:}{amber}{#1} {#2}\xspace}
\newcommand{\sdN}[1]{{{#1}}}
\newcommand{\sdNO}[1]{\red{ {#1}}} %{\ponders{SD:}{amber}{ } {#1}\xspace}
\newcommand{\julian}[1]{{#1}\xspace}
\newcommand{\sdM}[1]{#1\xspace}

\newcommand{\sophiaPonder}[2][]{\ponders{Sophia}{blue}{#1}{#2}\xspace}
\renewcommand{\sophia}[2][]{}

\newcommand{\sdfootnote}[1]{}



\begin{document}

 \author{Julian Mackay}
%\authornote{with author2 note}          %% \authornote is optional;
                                        %% can be repeated if necessary
\orcid{0000-0003-3098-3901}             %% \orcid is optional
\affiliation{
  %\position{Position2a}
  %\department{Engineering and Computer Science}             %% \department is recommended
  \institution{Victoria University of Wellington}           %% \institution is required
  %\streetaddress{Street2a Address2a}
  %\city{City2a}
  %\state{State2a}
  %\postcode{Post-Code2a}
   \country{NZ}   % \country{New Zealand}                   %% \country is recommended
}
\email{julian.mackay@ecs.vuw.ac.nz}         %% \email is recommended


\author{Susan Eisenbach}
%\authornote{with author2 note}          %% \authornote is optional;
                                        %% can be repeated if necessary
\orcid{0000-0001-9072-6689}             %% \orcid is optional
\affiliation{
  %\position{Position2a}
  %\department{}             %% \department is recommended
  \institution{Imperial College London}           %% \institution is required
  %\streetaddress{Street2a Address2a}
  %\city{City2a}
  %\state{State2a}
  %\postcode{Post-Code2a}
  \country{UK} %   \country{United Kingdom}                   %% \country is recommended
}
\email{susan@imperial.ac.uk}         %% \email is recommended

\author{James Noble}
%\authornote{with author2 note}          %% \authornote is optional;
                                        %% can be repeated if necessary
\orcid{0000-0001-9036-5692}             %% \orcid is optional
\affiliation{
  %\position{Position2a}
  %\department{Department2a}             %% \department is recommended
  \institution{Creative Research \& Programming}           %% \institution is required
 % \streetaddress{5 Fernlea Ave, Darkest Karori}
  \city{Wellington}
  %\state{State2a}
  \postcode{6012}
   \country{NZ}   %  \country{New Zealand}                   %% \country is recommended
}
\email{kjx@acm.org}         %% \email is recommended

%% Author with single affiliation.
\author{Sophia Drossopoulou}
%\authornote{with author1 note}          %% \authornote is optional;
                                        %% can be repeated if necessary
\orcid{0000-0002-1993-1142}             %% \orcid is optional
\affiliation{
  %\position{Position1}
  %\department{Department1}              %% \department is recommended
  \institution{Imperial College London}            %% \institution is required
  %\streetaddress{Street1 Address1}
  %\city{City1}
  %\state{State1}
  %\postcode{Post-Code1}
  \country{UK} %  \country{United Kingdom}                    %% \country is recommended
}
\email{scd@imperial.ac.uk}          %% \email is recommended


%% Abstract
%% Note: \begin{abstract}...\end{abstract} environment must come
%% before \maketitle command

\newcommand{\GoodWallet}{\ensuremath{\mathcal G}dWallet}  
\newcommand{\GoodParty}{\ensuremath{\mathcal G}dParty}  
\newcommand{\buyer}{\ensuremath{\texttt{buyer}}}
\newcommand{\seller}{\ensuremath{\texttt{seller}}}
\newcommand{\price}{\ensuremath{\texttt{price}}}
\newcommand{\weight}{\ensuremath{\texttt{weight}}}

\newcommand{\Unaffected}{\ensuremath{{U\!na{f\!f}\!ected}}}
\newcommand{\GoodEscrow}{\ensuremath{\mathcal G}dEscrow} %{\mathcal E}scrow}  
\renewcommand{\obeys}[2]{#2(#1)}

\newcommand{\kwd}[1]{\textbf{#1}}

%% 2012 ACM Computing Classification System (CSS) concepts
%% Generate at 'http://dl.acm.org/ccs/ccs.cfm'.
\begin{CCSXML}
<ccs2012>
   <concept>
       <concept_id>10011007.10010940.10010992.10010993.10011683</concept_id>
       <concept_desc>Software and its engineering~Access protection</concept_desc>
       <concept_significance>300</concept_significance>
       </concept>
   <concept>
       <concept_id>10011007.10011074.10011099.10011692</concept_id>
       <concept_desc>Software and its engineering~Formal software verification</concept_desc>
       <concept_significance>500</concept_significance>
       </concept>
   <concept>
       <concept_id>10003752.10003790.10011741</concept_id>
       <concept_desc>Theory of computation~Hoare logic</concept_desc>
       <concept_significance>500</concept_significance>
       </concept>
   <concept>
       <concept_id>10011007.10011006.10011008.10011009.10011011</concept_id>
       <concept_desc>Software and its engineering~Object oriented languages</concept_desc>
       <concept_significance>300</concept_significance>
       </concept>
 </ccs2012>
\end{CCSXML}

\ccsdesc[300]{Software and its engineering~Access protection}
\ccsdesc[500]{Software and its engineering~Formal software verification}
\ccsdesc[500]{Theory of computation~Hoare logic}
\ccsdesc[300]{Object oriented programming~Object capabilities}

%% End of generated code
 
\title{Safe Mediation in the Presence of Malicious Agents}
%% Keywords
%% comma separated list
%%%%%%%%%%%%%%%%%%%%%\keywords{keyword1, keyword2, keyword3}  %% \keywords are mandatory in final camera-ready submission
 
%% Note: \maketitle command must come after title commands, author
%% commands, abstract environment, Computing Classification System
%% environment and commands, and keywords command.


 % \input{abstract} % must come before maetitle

\begin{abstract}
xxx
xxx
xxx
\end{abstract}


\maketitle 




\section{Introduction}

Open world, where external, unknown, untrusted objects interact with internal, trusted objects, is upon us.
.. say who has worked on this ...
These works address the question of which program properties are preserved across external call boundaries, and 
establish guarantees that hold unconditionally — even in the presence of a malicious callee.

But they do not address the problem of \emph{mediation}, where  an object/function mediates a contract/transaction? between 
several other participant objects (agents? participants) each of which may, or may not be adhering to a protocol. 
The mediator has  the remit to ensure that if all participant objects indeed adhere to  protocol expected from them,  then the transaction 
will be implemented, 
but if one participant object does not adhere to the protocol, then none of the objects which adhere should be affected.
And in addition, the mediator has to be able to do all that in the absence of any central trusted authority.
In other words, the mediator has no way to find out whether a particular participant object adheres to the expected protocol.

%$\kwd{specification}$ 
 \begin{figure*}[t]
\begin{lstlisting}[mathescape=true, language=Chainmail, frame=lines]
$\GoodEscrow(\texttt{e}) \triangleq$
    $\kwd{exports}$
        $\kwd{method}$ deal($\buyer$:Object,$\seller$:Object,$\price$:Nat,$\weight$:Nat): bool
    $\kwd{such that}$     
        $true$
           $\kwd{\{}$ res := e.exchange($\buyer$,$\seller$,$\price$,$\weight$) $\kwd\}$
        res $\Longrightarrow$
            // 1$^{st}$ case:
            $\obeys  {\buyer} {\GoodParty}\ \wedge\ \obeys  {\seller} {\GoodParty}\ \wedge\   Enough\_Money\_Goods  \ \wedge$
            $Exchange\_Completed\  \wedge\  \Unaffected(others)$
            $\vee$
            // 2$^{nd}$ case:
            $\neg \obeys  {\buyer} {\GoodParty}\ \wedge\ \neg \obeys  {\seller} {\GoodParty}\ \wedge$
            $\Unaffected(protected\_others)$
        $\wedge$
        $\neg$ res $\Longrightarrow$
            // 3$^{rd}$ case:
            ($\, \obeys  {\buyer} {\GoodParty}\, \wedge\, \neg \obeys  {\seller} {\GoodParty}\ \vee\ ...   \mbox{similar cases} ...  \ )\  \ \wedge$   
            $\Unaffected(protected\_others)$
            $\vee$
            // 4$^{th}$ case:
            $\neg \obeys  {\buyer} {\GoodParty}\ \wedge\ \neg \obeys  {\seller} {\GoodParty} \ \wedge$
            $\Unaffected(protected\_others)$
            
\end{lstlisting}
\caption{Partial specification of $\GoodEscrow$.   The predicate $\GoodParty$ is defined in Fig. 2 and the meanings of $\mathit{Enough\_Money\_Goods}$ and $\mathit{Exchange\_Completed}$ should be self-evident. The definitions of $\mathit{Unaffected}$, $\mathit{others}$, and $\mathit{protected\_others}$ are omitted for now — the intuition should be clear, but their full definitions will require some care.}
\label{fig:EscrowSpec}
 \end{figure*}


In order to make it possible to implement such mediators, the participant objects   provide means for {\emph{attestation}}:
an object $o$ expresses that it believes that another object $o'$ adheres to protocol. 
Such an attestation however is {\emph{hypothetical:}} If $o$ itself is malicious, then the attestation it provided is not trustworthy, and $o'$
may or may not adhere to the protocol. 
If $o$ does indeed adhere to protocol, and makes an attestation for $o'$, then indeed $o'$ adheres to protocol.


Such questions have been studied by Marks Miller in ... and by ... Enumerate the examples he looked at .... First stab at reasoning in 
worhhop paper.

... outline the code of Escrow?  or James's beautiful diagram and write text to explain?



The specification of the Escrow, therefore would look as  in  Fig. \ref{fig:EscrowSpec}. ... 
Explain what all this means. ... 
Note that the 2nd and 4th case happen under the same conditions
In the end, the escrow does not know ... but if a participant ...

\newpage


As we said earlier, for escrow to be able to provide this guarantees, the participants need to provide a way of attestation, This is shown  in Fig. \ref{fig:GoodWallet}.

 \begin{figure*}[t]
\begin{lstlisting}[mathescape=true, language=Chainmail, frame=lines]
 $\GoodParty(\texttt{p}) \triangleq$
    $\kwd{exports}$
        $\kwd{field}$ money, goods: Object
    $\kwd{such that}$  
         $\obeys {\texttt{p.money}} {\GoodWallet}\ \wedge\ \obeys {\texttt{p.goods}} {\GoodWallet}$   
    
$\GoodWallet(\texttt{w}) \triangleq$
    $\kwd{exports}$
         $\kwd{method}$ transfer(w':Object,amt:Nat): bool
         $\kwd{method}$ make( ): Object
         $\kwd{ghost}$ balance: Nat 
         $\kwd{ghost}$ CanTrade(x): bool      
    $\kwd{such that}$     
        $\forall $ w'.$[  $ w.CanTrade(w') $\triangleq$ w.transfer(w',0)= $true \ ]$
        $\forall $ w'.$[  $ w.CanTrade(w')  $\Longrightarrow \ \obeys  {\texttt{w'}} {$\GoodWallet$} \ ]$
        $\forall $ w'.$[ $ w.make()=w' $ \Longrightarrow$ w.CanTrade(w') $\wedge$ w.balance=0 $]$
 
       $true$
          $\kwd{\{}$ res:= w.transfer(w',amt) $\kwd\}$
       res $\Longrightarrow$ w.CanTrade(w') $\wedge\ Enough\_Money\ \wedge\ Money\_Transfered \wedge\, \Unaffected(others)$
       $\wedge$
       $\neg$ res $\Longrightarrow (\neg$w.CanTrade(w')$\vee\, \neg Enough\_Money)\, \wedge\, \Unaffected(all)$
     
}
\end{lstlisting}
\caption{Partial specification of $\mathit{GoodParty}$ and $\mathit{GoodWallet}$. The meanings of $\mathit{Enough_Money}$ and $\mathit{Money_Transferred}$ should be self-evident. The definitions of $\mathit{Unaffected}$, $\mathit{others}$, and $\mathit{all}$ are simlar to to but slightly different than in Fig. 1 -- while their intended meanings should be clear from context, a precise treatment will reveal some non-trivial subtleties.}
\label{fig:GoodWallet}
 \end{figure*} 
 
 
Notice that our predicates have the  interesting characteristics, that  they talk about the hypothetical execution of method calls -- eg $\GoodWallet$ talks about the execution of a method call . Also, they   are recursively defined -- eg $\GoodWallet(gp)$  is defined among others, in terms of $\GoodWallet(gp.\texttt{sprout()})$. These questions make a well-founded definition more challenging.

In this paper we 

\begin{itemize}
\item
We propose a specification language to describe predicates like $\GoodWallet$ and $Escrow$
\item
We give semantics to these predicates, addressing the hypothetical method calls and the recursive definitions
\item
Develop a Hoare Logic to reason about such specifications, and deal with the hypothetical nature of protocol adgerence
\item
Outline a proof that a particular class sarisfies $GoodEscriw$, and another two strisfy $\GoodWallet$
\item
Outline how reasoning about our specifications can be mapped in a meaning-preserving way to a classical program verifier, and verify EScrow in Dafny

\end{itemize}

Here we say how we differ from our earlier paper.

%\section{Examples}  
%\label{s:NewExamples}
%\renewcommand{\this}{\texttt{this}}
\renewcommand{\balance}{\texttt{balance}}
 
  \begin{figure*}[t]
\begin{lstlisting}[mathescape=true, language=Chainmail, frame=lines]
$\textbf{specification}$ $GoodPurse(\this)$ {
    // $\mbox{we implicitly have that\ } \obeys {\this} {GoodPurse}$
    
    $\textbf{ghost}$ balance: Nat
     
    $\textbf{ghost}$ CanTrade(x) : bool
    
    $\textbf{method}$ sprout(): GoodPurse
    $\textbf{method}$ transfer($p'$:Any,$amt$:Nat) : bool
     
    $\forall p.[ \ \this.CanTrade(p)\  \triangleq\ \this.\texttt{transfer}(p,0) = \texttt{true} \ ]$
     
    // CanTrade $\mbox{implies}$ GoodPurse
     $\forall  p. [\ \this.CanTrade(p) \ \Longrightarrow \ \obeys  {p} {GoodPurse} \ ]$
       
    // CanTrade $\mbox{is an equivalence relation}$
    $\this.CanTrade(\this)$
    $\forall  p. [\ \this.CanTrade(p) \ \Longrightarrow \ p.CanTrade(\this) \ ]$
    $\forall  p,p' [\ \this.CanTrade(p) \wedge \ p.CanTrade(p')  \ \Longrightarrow \ \this.CanTrade(p)  \ ]$
       
    $\forall gp, [\ \texttt{this.sprout()}=gp \ \Longrightarrow  \ \texttt{this.CanTrade}(gp)\  \wedge\  p.\balance=0]$
    
    $\textbf{scoped-invr}\  \forall \prg b:\prg{int}.[\ \inside{\prg{this}}  \ \wedge \prg{this}.\prg{balance}=\prg b \ ]$ 
    
    // $CanTrade \ \mbox{is immutable and so is} \neg CanTrade$
    $\textbf{scoped-invr}\  \forall \prg p.[\ \ \prg{this}.CanTrade(p) \ ]$ 
    $\textbf{scoped-invr}\  \forall \prg p.[\ \ \neg \prg{this}.CanTrade(p) \ ]$ 
    
    // QUESTION: Is it implicit that $\obeys x {GoodPurse}$, \ $\neg (\obeys x {GoodPurse})$  are immutable?
         
    ${\texttt {this}}.CanTrade(p) \ \wedge\  {\texttt {amt}}\leq \texttt {this.balance}\ \wedge\ p\neq \this$
        $\textbf{\{}$ res:=this.transfer($p$,$amt$) $\textbf\}$
    res $\wedge$
    $\mbox{the balances  of \texttt{this} and } p \mbox{ change as should}$
    $\mbox{the balances of all other } GoodPurse \mbox{'s remain unaffected}$
     
    $\neg {\texttt {this}}.CanTrade(p) \vee {\texttt {amt}}> \texttt {this.balance}$
        $\textbf{\{}$ res:=this.transfer($p$,$amt$) $\textbf\}$
    $\neg$ res $\wedge$
     $\mbox{the balances of all } GoodPurse \mbox{'s remain unaffected}$
     
}
\end{lstlisting}
\caption{Specification of  $GoodPurse$ }
\label{fig:GoodPurse}
 \end{figure*}
 
 \vspace{.2cm}
 
 \textbf{Notes on} $GoodPurse$:\\
 \begin{itemize}
 \item
 \texttt{CanTrade} is not ghost
 \item
 Should we have have written  \\
$\strut \hspace{1cm} \textbf{scoped-invr}\  \forall \prg b:\prg{int}.\forall p.[\ \obeys p {GoodPurse} \ \wedge\ \inside{\prg{p}}  \ \wedge \prg{p}.\prg{balance}=\prg b \ ]$ 
\\
instead of\\
$\strut \hspace{1cm} \textbf{scoped-invr}\  \forall \prg b:\prg{int}.[\ \inside{\prg{this}}  \ \wedge \prg{this}.\prg{balance}=\prg b \ ]$ 
\\
and similar for all appearances of \texttt{this}? 
\end{itemize}  
 
 
 \begin{figure*}[t]
\begin{lstlisting}[mathescape=true, language=Chainmail, frame=lines]
$\textbf{specification}$ $FairParty$ {
    
     $\textbf{field}$  moneyPurse 
     $\textbf{field}$  foodPurse 
     
     $\textbf{ghost}$ CanExchng(x) :       
     
     $\forall fp. [\ \obeys  {fp} {FairParty} \  \ \Longrightarrow \ \obeys  {fp.\prg{moneyPurse} }  {GoodPurse} \ \wedge\ \obeys  {fp.\prg{foodPurse} }  {GoodPurse} \ ]$
      
     $\forall fp. \texttt{this}.CanExchng(fp) \ \triangleq \texttt{this.moneyPurse}.CanTrade(\texttt{fp.moneyPurse}) \wedge$   
                          $\texttt{this.foodPurse}.CanTrade(\texttt{fp.foodPurse}) $

}
\end{lstlisting}
\caption{Specification of  $FairParty$}
\label{fig:FairParty}
 \end{figure*}
 
  \textbf{Notes on} $FairParty$:\\
 \begin{itemize}
 \item
 \texttt{CanExchng} is ghost, as opposed to \texttt{CanTrade}\\
\item
Should we have defined\\
$\strut \hspace{1cm} \textbf{field}$  \texttt{moneyPurse: GoodPurse};  $\textbf{field}$;  \texttt{moneyPurse: GoodPurse}. \\
If we do that, then we could omit line 8 from Fig. \ref{fig:FairParty}
\item
A corollary of the specification is that \\
$\strut \hspace{1cm} \forall p. [\ \texttt{this}.CanExchng(p) \Longrightarrow \ \obeys p {FairParty}\ ]$
 \end{itemize}  
 
% \footnote{Should we have defined instead 
%\begin{lstlisting}[mathescape=true, language=Chainmail, frame=lines]
%$\textbf{specification}$ $FairParty$ {
%    
%     $\textbf{field}$  moneyPurse : $GoodPurse$
%     $\textbf{field}$  foodPurse : : $GoodPurse$
%     
%     $\textbf{ghost}$ CanExchng(x) :       
%          
%      $\forall x. \texttt{this}.CanExchng(x) \ \triangleq \texttt{this.moneyPurse}.CanTrade(\texttt{this.moneyPurse}) \wedge$   
%                          $\texttt{this.foodPurse}.CanTrade(\texttt{this.foodPurse}) $
%
%}
%\end{lstlisting}
%}

 \begin{figure*}[t]
\begin{lstlisting}[mathescape=true, language=Chainmail, frame=lines]
$\textbf{specification}$ $Escrow$ {
    
    // 1$^{st}$ case:
    $\obeys  {buyer} {FairParty}$ $\wedge$ 
    $buyer$.CanExchng($seller$)  $\wedge  \mbox{ enough money and goods }$  
        $\textbf{\{}$ res:=this.deal($buyer$,$seller$,$price$,$weight$) $\textbf\}$
    res $\wedge$
   $\mbox{ exchange happened } \wedge  \mbox{ all other moneys and goods remain unaffected }$
   
    //  2$^{nd}$ case:
    $\{ p, p' \} = \{$ buyer, seller $\} \ \wedge\  \obeys  {p} {FairParty}$ $\wedge$ 
    $p$.CanExchng($p'$) $\wedge \ \neg(\mbox{ enough money and goods})$  
        $\textbf{\{}$ res:=this.deal($buyer$,$seller$,$price$,$weight$) $\textbf\}$
    $\neg$res $\wedge$
    $\mbox{ exchange did not happen}  \wedge  \mbox{ all other moneys and goods remain unaffected }$
    
    //  3$^{rd}$ case:
    $\{ p, p' \} = \{$ buyer, seller $\} \ \wedge\  \obeys  {p} {FairParty}$ $\wedge$ 
    $\neg p$.CanExchng($p'$) 
        $\textbf{\{}$ res:=this.deal($buyer$,$seller$,$price$,$weight$) $\textbf\}$
    $\neg$res $\wedge$
    $\mbox{ moneys and goods protected from } p' \mbox{ remain unaffected }$   
    
    //  4$^{th}$ case:
    $\neg \obeys  {buyer} {FairParty}$ $\wedge$ $\neg \obeys  {seller} {FairParty}$ 
        $\textbf{\{}$ res:=this.deal($buyer$,$seller$,$price$,$weight$) $\textbf\}$
    res$\in\{$ true, false $\}$ $\wedge$
   $\mbox{ moneys and goods protected from } seller  \mbox{ and } buyer \mbox{ remain unaffected }$   
}

\end{lstlisting}
\caption{Specification of  $Escrow$.\prg{::deal} }
\label{fig:EscrowSpec}
 \end{figure*}
 
 
  \textbf{Notes on} $FairParty$:\\
 \begin{itemize}
 \item
The 2$^{nd}$ case from above can is equivalent to\\
\begin{lstlisting}[mathescape=true, language=Chainmail, frame=lines]
....   
    //  2a  case:
    $\obeys  {buyer} {FairParty}$ $\wedge$ 
    $buyer$.CanExchng($seller$) $\wedge  \neg(   \mbox{ enough money and goods })$  
        $\textbf{\{}$ res:=this.deal($buyer$,$seller$,$price$,$weight$) $\textbf\}$
    $\neg$res $\wedge$
    $\mbox{ exchange did not happen}  \wedge  \mbox{ all other moneys and goods remain unaffected }$
    
    //  2b case:
    $\obeys  {seller} {FairParty}$ $\wedge$ 
    $seller$.CanExchng($buyer$) $\wedge  \neg(   \mbox{ enough money and goods })$  
        $\textbf{\{}$ res:=this.deal($buyer$,$seller$,$price$,$weight$) $\textbf\}$
    $\neg$ res $\wedge$
    $\mbox{ exchange did not happen}  \wedge  \mbox{ all other moneys and goods remain unaffected }$
\end{lstlisting}

\item
Similar expandion is possible for the 3$^{rd}$ case.
\item
A corollary of the specification is that \\
$\strut \hspace{1cm} \forall p. [\ \texttt{this}.CanExchng(p) \Longrightarrow \ \obeys p {FairParty}\ ]$
 \end{itemize}  
 
 
\newpage

Here the longer version, before Susan's clever suggestion

 \begin{lstlisting}[mathescape=true, language=Chainmail, frame=lines]
$\textbf{specification}$ $Escrow$ {
    
    // 1$^{st}$ case:
    $\obeys  {buyer} {FairParty}$ $\wedge$ 
    $buyer$.CanExchng($seller$)  $\wedge  \mbox{ enough money and goods }$  
               $\textbf{\{}$ res:=this.deal($buyer$,$seller$,$price$,$weight$) $\textbf\}$
    res $\wedge$
   $\mbox{ exchange happened } \wedge  \mbox{ all other moneys and goods remain unaffected }$
   
    //  2$^{nd}$ case:
    $\obeys  {buyer} {FairParty}$ $\wedge$ 
    $buyer$.CanExchng($seller$) $\wedge  \neg(   \mbox{ enough money and goods })$  
               $\textbf{\{}$ res:=this.deal($buyer$,$seller$,$price$,$weight$) $\textbf\}$
    $\neg$res $\wedge$
    $\mbox{ exchange did not happen}  \wedge  \mbox{ all other moneys and goods remain unaffected }$
    
    //  3$^{rd}$ case:
    $\obeys  {seller} {FairParty}$ $\wedge$ 
    $seller$.CanExchng($buyer$) $\wedge  \neg(   \mbox{ enough money and goods })$  
               $\textbf{\{}$ res:=this.deal($buyer$,$seller$,$price$,$weight$) $\textbf\}$
    $\neg$ res $\wedge$
    $\mbox{ exchange did not happen}  \wedge  \mbox{ all other moneys and goods remain unaffected }$
    
     //  4$^{th}$ case:
    $\obeys  {buyer} {FairParty}$ $\wedge$ 
    $\neg (buyer$.CanExchng($seller$))  
               $\textbf{\{}$ res:=this.deal($buyer$,$seller$,$price$,$weight$) $\textbf\}$
    $\neg$res $\wedge$
    $\mbox{ moneys and goods protected from } seller \mbox{ remain unaffected }$   
    
       //  5$^{th}$ case:
    $\obeys  {seller} {FairParty}$ $\wedge$ 
    $\neg (seller$.CanExchng($buyer$))  
               $\textbf{\{}$ res:=this.deal($buyer$,$seller$,$price$,$weight$) $\textbf\}$
    $\neg$ es $\wedge$
   $\mbox{ moneys and goods protected from } buyer \mbox{ remain unaffected }$   
 
    //  6$^{th}$ case:
    $\neg \obeys  {buyer} {FairParty}$ $\wedge$   $\neg \obeys  {seller} {FairParty}$ 
               $\textbf{\{}$ res:=this.deal($buyer$,$seller$,$price$,$weight$) $\textbf\}$
    res$\in\{$ true, false $\}$ $\wedge$
   $\mbox{ moneys and goods protected from } seller  \mbox{ and } buyer \mbox{ remain unaffected }$   
}
\end{lstlisting}

 

\end{document}
